
\newcommand{\SKIP}[1]{}

%-------------constructors--------------------

\newcommand{\tld}[1]{\tilde{#1}}


\newcommand{\nil}{[]}
\newcommand{\cons}[2]{#1::#2}
\newcommand{\der}[2]{#1\hat{\,\,}#2}
\newcommand{\cut}[2]{\mathsf{cut}(#1,#2)}
\newcommand{\tcut}[2]{\mathsf{cut}_\tau(#1,#2)}  %----- tight cut


%\newcommand{\mut}{\tilde{\mu}}%------------ mu tilde

%\newcommand{\lt}[3]{\mathsf{let}\,#2\,\mathsf{be}\,#1\,\mathsf{in}\,#3}%------- let expression
%\newcommand{\lt}[3]{\mathsf{let}\,#2:=#1\,\mathsf{in}\,#3}%------- let expression
\newcommand{\hd}[1]{\mathsf{hd}(#1)} %-------- coercion
\newcommand{\thd}[1]{\mathsf{hd}_\tau(#1)} %-------- tight coercion
\newcommand{\app}[1]{\mathsf{app}(#1)} %------------ coercion
\newcommand{\tapp}[1]{\mathsf{app}_\tau(#1)} %------------ tight coercion

%----------- substitution ----------
\newcommand{\sub}[3]{[#1/#2]#3} %----------- substitution
\newcommand{\tsub}[3]{[#1/_{\!\!\tau}#2]#3} %----------- tight substitution

%----------------- contexts ----------------
\newcommand{\hol}{[\cdot]} %--- hole
\newcommand{\fh}[1]{[#1]} %------ full hole

% \newcommand{\HH}{\mathcal H}
% \newcommand{\K}{\mathcal K}
% \newcommand{\C}{\mathcal C}
% \newcommand{\D}{\mathcal D}



%---------explicit substitutions----------
\newcommand{\la}{\langle}
\newcommand{\ra}{\rangle}
\newcommand{\x}{\mathtt{x}}
%\newcommand{\xs}[3]{\la#1/#2\ra#3}

\newcommand{\xsub}[3]{\la #1/#2\ra #3} %--------explicit substitution
\newcommand{\mcut}[3]{\{#1=:#2\} #3} %--------mid-cut

%--------- Reduction rules names---------------
\newcommand{\tbeta}{\beta_\tau} %------ tight beta

%-----------lambdas-----------------------
\newcommand{\lb}{\lambda}
%\newcommand{\lm}{\lambda_m}
%\newcommand{\lx}{\lambda \x}
%\newcommand{\lmu}{\lambda\mu}

\newcommand{\lbd}{\boldsymbol{\lb}} %------ bold lambda to denote the ordinary lambda-calc

\newcommand{\olb}{\overline{\lambda}}
\newcommand{\tolb}{\overrightarrow{\lb}} %-------- tight \olb
\newcommand{\ulb}{\underline{\lambda}}
\newcommand{\tulb}{\underleftarrow{\lb}} %----- tight \ulb



\newcommand{\nvec}{\underleftarrow{\,\,\lb}} %------ lambda vec in natural deduction
\newcommand{\rnvec}{\stackrel{\large\lb}{\tiny\hookleftarrow}} %-- relaxed lambda vec in natural deduction

%\newcommand{\lmmt}{\olb\mu\mut}


%-------- sets of expressions
\newcommand{\VTerms}{VecTerms} %----------- vectorial terms
\newcommand{\tVTerms}{VecTerms_\tau} %----------- tight vectorial terms
\newcommand{\Lists}{Lists}
\newcommand{\tLists}{Lists_\tau}
\newcommand{\BTerms}{BidTerms} %------------ bidirectional terms
\newcommand{\tBTerms}{BidTerms_\tau} %------------ tight bidirectional terms
\newcommand{\Elims}{Elims}
\newcommand{\tElims}{Elims_\tau}

%---------------maps and their symbols-------------------

\newcommand{\fgt}[1]{|#1|} %---------------- forgetful map
\newcommand{\bid}[1]{#1^{\natural}}  %---- from lambda to bid-lambda

%------------ formulas/types
\newcommand{\imp}{\supset} %--- implication

%------------ logics
\newcommand{\LJT}{LJT}
\newcommand{\NJT}{NJT}


%---------- sequents for intuitionistic logic
\newcommand{\seqt}[2]{#1\vdash#2} %------
\newcommand{\Seqt}[3]{#1\vdash#2:#3} %---- term sequent

\newcommand{\seql}[3]{#1;#2\vdash#3} %------
\newcommand{\Seql}[4]{#1;#3:#2\vdash #4} %---- list sequent

\newcommand{\seqP}[2]{#1\vdash#2} %------
\newcommand{\SeqP}[3]{#1\vdash#2:#3} %---- term sequent (P's in nat ded)

\newcommand{\seqE}[2]{#1\rhd#2} %------
\newcommand{\SeqE}[3]{#1\rhd#2:#3} %----  sequent for eliminations (E's)

% --------------------
% lstlisting coq style (inspired from a file of Assia Mahboubi
\lstdefinelanguage{Coq}%
  {morekeywords={Variable,Inductive,CoInductive,Fixpoint,CoFixpoint,%
      Definition,Lemma,Theorem,Corollary,Axiom,Local,Save,Grammar,Syntax,intro,%
      trivial,Qed,intros,Symmetry,Simpl,Rewrite,apply,elim,assumption,%   
      left,Case,auto,unfold,exact,right,Hypothesis,Pattern,destruct,%     Cut
      constructor,Defined,Fix,Record,Proof,induction,Hints,exists,let,in,%
      Parameter,split,Red,reflexivity,transitivity,if,then,else,Opaque,%
      Transparent,inversion,Absurd,Generalize,Mutual,Cases,of,end,Analyze,%
      AutoRewrite,Functional,Scheme,params,Refine,using,discriminate,Try,%
      Require,Load,Import,Scope,Set,Open,Section,End,match,with,Ltac,%
      Instance,%
      bind,as,%
      % , exists, forall
	},%
   sensitive, %
   morecomment=[n]{(*}{*)},%
   morestring=[d]",%
   literate=
   {=>}{{$\Rightarrow$}}1
   {>->}{{$\rightarrowtail$}}2
   {<->}{{$\leftrightarrow$}}1
   {->}{{$\rightarrow$}}1
   {\/\\}{{$\wedge$}}1
   {|-}{{$\vdash$}}1
   {\\\/}{{$\vee$}}1
   {~}{{$\sim$}}1
   {exists}{{$\exists\!\!$}}1
   {forall}{{$\forall\!\!$}}1
   {beta1}{{$\beta_1$}}1
   {beta2}{{$\beta_2$}}1
   {pi}{{$\pi$}}1
   {epsilon}{{$\epsilon$}}1
   {piepsilon}{{$\pi\epsilon$}}2
   {sigma}{{$\sigma$}}1
   {theta}{{$\theta$}}1
   {tau}{{$\tau$}}1
   {rho}{{$\rho$}}1
   {xi}{{$\xi$}}1
   {iota}{{$\iota$}}1
   {Axiom}{Axiom}3 % patch
   {exi}{exi}3 % insufficient patch to print reflexivity correctly
   {zeta}{{$\zeta$}}1
   {Gamma}{{$\Gamma$}}1
   {Delta}{{$\Delta$}}1
   {\\rhd}{{$\rhd$}}1
   %{>>}{\scomp}1
   %{<>}{{$\neq$}}1 indeed... no.
  }[keywords,comments,strings]%
\lstset{
   basicstyle=\ttfamily,
   keywordstyle=\bfseries\color{blue}
  }
\lstset{language=Coq}
\lstset{columns=fullflexible, keepspaces}

%---------- inline codeblocks
\newcommand{\lst}{\lstinline}

%--- theoremstyles
% \theoremstyle{plain}
% \newtheorem{theorem}{Theorem}
% \newtheorem{proposition}{Proposition}
% \newtheorem{lemma}{Lemma}
% \newtheorem{corollary}{Corollary}

% \theoremstyle{definition}
% \newtheorem{definition}{Definition}

% \theoremstyle{remark}
% \newtheorem{remark}{Remark}

\providecommand{\theoremautorefname}{Theorem}
\providecommand{\propositionautorefname}{Proposition}
\providecommand{\lemmaautorefname}{Lemma}
\providecommand{\corollaryautorefname}{Corollary}

\newenvironment{autoproof}[1]
{\begin{proof}[Proof of \autoref{#1}]} {\end{proof}}

% ------- Rocq Links
\newcommand{\baselink}{https://w3.math.uminho.pt/~luis/code/html/}
\newcommand{\rocqlink}[1]{%
  \,\href{\baselink#1}{\raisebox{-0.20ex}{\includegraphics[height=0.9em]{icon-rocq-orange.png}}}%
}
\newcommand{\rocqlogo}{%
 {\raisebox{-0.20ex}{\includegraphics[height=0.9em]{icon-rocq-orange.png}}}%
}
